% ============================================================
% APPENDIX D: VALUE CANONICALIZATION
% ============================================================

\chapter{Canonical Values}
\label{app:values}

The Meridian framework predicts a one-parameter family of cosmologies indexed by $\zz = \xi\,\Phi_0^2/M_5^3$. Different chapters evaluate this family at different observational benchmarks, producing a range of numerical values for derived quantities---all consistent, but potentially confusing without a single reference table. This appendix resolves the proliferation by organising every derived quantity into three tiers according to its epistemic status.

\paragraph{Resolution principle.} The curve $w_0(\zz)$ is the prediction, not any single number. Values that appear to conflict (e.g.\ $w_0 = -0.830$ in Chapter~\ref{ch:ncg} vs.\ $w_0 = -0.990$ in Chapter~\ref{ch:observational}) are the \emph{same} parametric prediction evaluated at different $\zz$ benchmarks.

%% ====================================================================
\section{Tier~1: Structure-Determined}
\label{app:val-tier1}
%% ====================================================================

These values are derived from the geometry and algebra alone. They depend on neither $\zz$ nor $\eps_1$; they are exact consequences of the axioms A1--A6.

\begin{table}[h!]
\caption{Tier~1 quantities: structure-determined, zero free parameters.}
\label{tab:val-tier1}
\small
\renewcommand{\arraystretch}{1.25}
\begin{tabular}{@{}p{3.5cm}p{3.0cm}p{5.5cm}p{2.5cm}@{}}
\toprule
\textbf{Quantity} & \textbf{Value} & \textbf{Source} & \textbf{Confidence} \\
\midrule

$\xi$ (conformal coupling) &
$1/6$ &
Three complementary perspectives on a single geometric fact: Seeley--DeWitt $(a_2)$ coefficient, Weyl conformal-invariance of the kinetic term, and radion-as-metric-fluctuation (see App.~\ref{app:xi}) &
Very high \\

$C_\mathrm{GB}$ (Gauss--Bonnet coeff.) &
$2/3$ &
Exact; three independent derivations &
Very high \\

$N_g$ (generation number) &
$3$ &
Algebraic maximum: $J_3(\mathbb{O})$ pins $N_g \leq 3$; three independent perspectives (Dixon, anomaly cancellation, Clifford) confirm the maximum is saturated &
Very high \\

$\alpha_T = \alpha_B = \alpha_M$ &
$0$ (exactly) &
Cuscuton constraint + GR perturbation structure &
Very high \\

$w_a$ &
$\approx 0$ (structurally) &
No phantom crossing (topological barrier); $|w_a| \lesssim 0.02$ &
High \\

$\eta$ (gravitational slip) &
$1$ &
Cuscuton: no scalar propagation $\Rightarrow$ no anisotropic stress &
Very high \\

Growth index $\gamma$ &
$0.55$ (GR value) &
Growth--expansion decoupling: $\mu = \Sigma = 1$ &
High \\

\bottomrule
\end{tabular}
\end{table}

%% ====================================================================
\section{Tier~2: Parametric}
\label{app:val-tier2}
%% ====================================================================

These values depend on the GB correction parameter $\eps_1 = 0.010 \pm 0.002$, which is derived from the spectral action (Chapter~\ref{ch:ncg}) but whose precise value depends on the spectral cutoff function choice ($\pm 11\%$ across four cutoff functions).

\begin{table}[h!]
\caption{Tier~2 quantities: one-parameter family depending on $\eps_1$.}
\label{tab:val-tier2}
\small
\renewcommand{\arraystretch}{1.25}
\begin{tabular}{@{}p{3.0cm}p{3.0cm}p{5.5cm}p{2.5cm}@{}}
\toprule
\textbf{Quantity} & \textbf{Value} & \textbf{Formula / Source} & \textbf{Confidence} \\
\midrule

$\eps_1$ (GB correction) &
$0.010 \pm 0.002$ &
$d = 5$ Seeley--DeWitt with corrected Weyl decomposition; $C_\mathrm{GB} = 2/3$ &
High \\

$C_\mathrm{KK}$ &
$(1.64 \pm 0.33) \times 10^{-4}$ &
$\frac{(1 + q_0)^2\,\Omega_\mathrm{DE}\,\eps_1}{4(1 - q_0)^2}$ at Planck 2018 fiducial &
High \\

$c_s$ (sound speed) &
$[12c,\; 15c]$ &
$\sqrt{C_q/\eps_1}$; $C_q$ depends on deceleration parameter benchmark &
High \\

$c_s^2$ &
$\sim 216$ &
$C_q/\eps_1$ at canonical benchmark &
High \\

$n_s$ (spectral index) &
$0.965 \pm 0.003$ &
K\"ahler modulus attractor inflation ($\alpha = 1$) &
High \\

$r$ (tensor-to-scalar) &
$0.004 \pm 0.001$ &
Same inflationary model; $N_* \approx 57$ &
High \\

$m_\mathrm{rad}$ (radion mass) &
$\approx 120~\mathrm{GeV}$ &
Quantum stabilisation (99.7\% from NCG one-loop) &
Medium \\

$m_{ee}$ (0$\nu\beta\beta$ mass) &
$1.5$--$5~\mathrm{meV}$ &
$S_3$ breaking of democratic Majorana matrix &
Medium \\

\bottomrule
\end{tabular}
\end{table}

\paragraph{Sound speed resolution.}
The value $c_s = 11.3c$ appearing in earlier drafts used $\eps_1 = 0.018$ from the pre-correction $d = 5$ Weyl decomposition. The corrected value is $\eps_1 = 0.010$, giving $c_s \in [12c,\, 15c]$. The pure cuscuton limit ($\eps_1 \to 0$, $c_s \to \infty$) is always stated as context, not as a prediction.

\paragraph{Scalar boundary value resolution.}
The value $\Phi_0 = 0.477$ appearing in earlier drafts was reverse-engineered from an assumed $\zz = 0.038$ rather than derived from the Israel junction conditions. The correct junction-condition solution gives $\Phi_0 = 0.076067$ (a factor-of-six correction), with $\zz = \xi\,\Phi_0^2/M_5^3 = 9.64 \times 10^{-4}$ at the JC benchmark. See App.~\ref{app:selftuning} (Historical note on $\Phi_0$) for the full derivation. All Tier~1 and Tier~2 values in this appendix use the corrected $\Phi_0 = 0.076$.

\paragraph{Bayes factor resolution.}
The Savage--Dickey Bayes factor $B_{10} = 171\!:\!1$ appearing in earlier drafts was computed using the historical best-fit $\zz = 0.038$ on the Hubble--Kristian subset alone. With the corrected JC benchmark $\zz \approx 10^{-3}$, the posterior shifts toward $\zz = 0$ and the Bayes factor is substantially smaller; individual-probe detection significance is $\lesssim 0.3\sigma$, and a multi-probe combination is required. The historical value is retained in App.~\ref{app:hk} for transparency but is not current evidence for $\zz \neq 0$.

%% ====================================================================
\section{Tier~3: Data-Constrained}
\label{app:val-tier3}
%% ====================================================================

The brane separation parameter $\zz$ is the single free parameter. Different observational probes constrain it at different values; these are not contradictions but different measurements of the same quantity.

\begin{table}[h!]
\caption{Tier~3 quantities: data-constrained benchmarks for $\zz$.}
\label{tab:val-tier3}
\small
\renewcommand{\arraystretch}{1.25}
\begin{tabular}{@{}p{3.5cm}p{2.2cm}p{1.8cm}p{4.5cm}p{2.5cm}@{}}
\toprule
\textbf{Benchmark} & \textbf{$\zz$} & \textbf{$w_0$} & \textbf{Context} & \textbf{Where Used} \\
\midrule

Spectral chain &
$8.8 \times 10^{-4}$ &
$-0.830$ &
Self-consistent GW solution at $\varepsilon = 0.275$ (DESI-fitted; see Appendix~\ref{app:gw}) &
Ch.~\ref{ch:ncg}, \S\ref{sec:4-brane-parameter} \\

CAMB best-fit &
$0.013$ &
$-0.987$ &
Planck CMB + DESI Y1 BAO combined Boltzmann fit &
Ch.~\ref{ch:observational} \\

\textbf{Weighted mean} &
$\mathbf{0.016 \pm 0.002}$ &
$\mathbf{-0.990}$ &
Four-probe inverse-variance weighted (HK + $H(z)$ + CAMB + multi-probe); $\chi^2/\mathrm{dof} = 1.10$ &
\textbf{Canonical} \\

Multi-probe &
$0.020$ &
$-0.993$ &
DESI DR2 BAO + $f\sigma_8$ + Planck $H_0$ + HK &
Ch.~\ref{ch:observational} \\

HK--CMB &
$0.037$ &
$-0.996$ &
Hiramatsu--Kobayashi $\beta_\mathrm{HK}$ from Planck &
Ch.~\ref{ch:observational} \\

\bottomrule
\end{tabular}
\end{table}

\paragraph{Which value to use.}
For canonical single-number summaries, use the weighted mean: $\zz = 0.016 \pm 0.002$, $w_0 = -0.990$. When discussing the spectral chain specifically (Chapter~\ref{ch:ncg}), use the junction-condition benchmark and clearly label it as depending on the fitted parameter $\varepsilon = 0.275$.

%% ====================================================================
\section{The Goldberger--Wise Result}
\label{app:val-gw-note}
%% ====================================================================

The GW parameter fix (Appendix~\ref{app:gw}) established that $\varepsilon$ is \emph{genuinely external}---it cannot be derived from the spectral action. The spectral chain value $w_0 = -0.830$ therefore depends on a fitted parameter ($\varepsilon = 0.275$ from DESI), while the data-driven weighted mean $w_0 = -0.990$ does not. This distinction is not a weakness but a feature: it identifies the precise boundary between what the NCG spectral triple predicts (geometry, cuscuton kinetic structure, coupling constants) and what requires brane stabilisation physics (the modulus stabilisation scale). The theory remains one-parameter ($\zz$ free), as claimed throughout.

%% ====================================================================
\section{Structured Apophaticism: Theory-Scope vs.\ Access Remainder}
\label{app:val-apophatic}
%% ====================================================================

Throughout the monograph we have followed a three-register discipline: a \emph{prose} register stating claims in ordinary physicist language; a \emph{formal} register carrying the derivations, bounds, and numerical outputs; and an \emph{apophatic remainder} that names what the current apparatus cannot say.  This last register is Principle-aligned: the Coherence Principle predicts that any stream-local map will have a remainder, and that naming it is part of the map's honesty.  In a post-publication review we found that the remainder is not monolithic.  It splits cleanly into two classes, and the split is load-bearing for scope discipline.

\paragraph{Theory-scope remainder.}
What the current NCG apparatus, extended by the cuscuton and the Goldberger--Wise sector as presently computed, cannot derive:
\begin{itemize}\setlength{\itemsep}{0pt}
  \item the orbifold stabilisation scale $k y_c$ (equivalently the hierarchy between the UV brane and the IR brane);
  \item the Goldberger--Wise scaling dimension $\varepsilon$, hence the factor-of-3 gap between naive $\varepsilon = \sqrt{2/3} \approx 0.816$ and DESI-fitted $\varepsilon = 0.275$ (see App.~\ref{app:gw-implications}, candidate inspection-depth ceiling);
  \item the specific empirical brane tensions $T_\pm$ (self-tuning fixes them up to the stabilisation scale, not independently);
  \item the first-order vs.\ crossover status of the Randall--Sundrum stabilisation transition (a dedicated bounce-action computation on the cuscuton-backreacted GW potential has not been performed);
  \item the exact normal-hierarchy Majorana phase structure that pins P5's $m_{ee}$ range to meV precision rather than percent precision.
\end{itemize}
These are residuals that a more careful computation \emph{within the current theoretical apparatus} might close.  Theory-scope remainder is accessible in principle through harder work inside the same framework.

\paragraph{Access remainder.}
Depth that is \emph{in principle} accessible through higher-dimensional or cross-framework apparatus, but lies outside the currently-published record available to this monograph:
\begin{itemize}\setlength{\itemsep}{0pt}
  \item the placement of our specific physical stream within the Ruliad or the string-landscape ensemble---which basin we occupy, and why, among the set of basins compatible with the same generic-physical-aspect invariants ($\xi = 1/6$, $\CGB = 2/3$, $N_g \leq 3$);
  \item the relationship between the NCG spectral triple and potentially-equivalent higher-categorical or non-commutative-geometric reformulations that might embed the Meridian structure into a larger family of streams;
  \item any structural consequences of inter-basin accessibility at the radion-mass scale (see App.~\ref{app:radion-thinness} for the discipline-controlled structural observation).
\end{itemize}
Access remainder is accessible in principle but not through the current monograph's apparatus alone; it would require importing results from a higher framework (Ruliad, string landscape, or cross-stream comparison literature).

\paragraph{Why the split matters.}
Collapsing the two classes into a single undifferentiated ``open questions'' list obscures a scope-control distinction that the Coherence Principle predicts should be preserved: a coherent description should be able to name which parts of its remainder are in-theory-residual and which are out-of-theory-access.  The former invites harder computation; the latter invites inter-framework dialogue.  The former is a falsifiability hook (if the theory-scope remainder is closed and still disagrees with data, the theory fails); the latter is a scope marker (if the access remainder is bridged by a higher apparatus, the theory is embedded rather than replaced).  Both are load-bearing; neither is evasive.  The monograph's claims are pitched strictly inside what the apparatus can presently say, with the remainder named in both classes and neither absorbed into the other.

%% ====================================================================
\section{Stream Scope Audit}
\label{app:val-stream-scope}
%% ====================================================================

A companion to the apophatic-remainder split (\S\ref{app:val-apophatic}) is a forward audit: for every derived quantity in the monograph, label whether it is a property of the \emph{physical aspect of reality as such} (generic-physical-aspect), a property of \emph{our specific stream/basin} within that aspect (this-stream), or a quantity whose scope-class is genuinely ambiguous and requires further work (saturation/ambiguous).  This is the map-of-what-kind-of-map-the-map-is.  The Principle predicts that a coherent description should be able to name its own scope at every derived quantity; v1 did not, v2 does.  The audit labels what is already derived; it does not relabel or rederive anything.

\paragraph{Generic-physical-aspect quantities.}
True of any physical stream that uses the same algebra, spectral triple, or AdS$_5$ geometry:
\begin{itemize}\setlength{\itemsep}{0pt}
  \item $\xi = 1/6$ (conformal coupling) --- follows from Seeley--DeWitt $a_2$, Weyl invariance, AdS$_5$ conformal geometry; any stream with the same scalar-as-metric-fluctuation identification gets the same $\xi$.
  \item $\CGB = 2/3$ (Gauss--Bonnet junction factor) --- fixed identically by the Davis junction conditions with GB term; no stream-specificity.
  \item $N_g \leq 3$ (generations bound) --- the algebraic maximum from $J_3(\mathbb{O})$; any stream using the octonionic exceptional Jordan algebra is capped at three.
  \item Existence of a cuscuton-like coherence-maintenance structure (zero-DOF first-order scalar) --- follows from the Horndeski uniqueness argument, not from any stream-specific choice.
  \item The absence of phantom crossing ($w \geq -1$ topologically) --- follows from the cuscuton kinetic-function positivity, a structural consequence of the square-root kinetic term.
\end{itemize}

\paragraph{This-stream quantities.}
Identify \emph{our} basin within the physical aspect:
\begin{itemize}\setlength{\itemsep}{0pt}
  \item $\zz$ (cuscuton mass-scale parameter, weighted mean $0.016 \pm 0.002$) --- fitted to our observed $w_0$.
  \item $\varepsilon = 0.275$ (Goldberger--Wise scaling dimension) --- fitted to DESI; the naive conformal-coupling value $\sqrt{2/3}$ corresponds to a different stream's stabilisation configuration (or to the same stream's generic-aspect expectation, with the fitted value carrying the stream-specific information).
  \item The specific numerical Yukawa couplings $y_{u,c,t}$, $y_{d,s,b}$, $y_{e,\mu,\tau}$ --- the octonionic spectral triple supplies the hierarchy structure; the overall normalisation is this-stream.
  \item $\CKK = (1.64 \pm 0.33)\times 10^{-4}$ --- carries a specific empirical input from the KK-mode decomposition.
  \item $m_\mathrm{rad} \approx 120~\mathrm{GeV}$ --- depends on specific GW-backreaction on the cuscuton (and see App.~\ref{app:radion-thinness}).
  \item $m_{ee}$ range at meV precision --- depends on the stream's Majorana phase structure.
\end{itemize}

\paragraph{Saturation / ambiguous quantities.}
May be generic or stream-specific; scope-class requires further work:
\begin{itemize}\setlength{\itemsep}{0pt}
  \item $N_g = 3$ (bound saturated, not just bounded) --- whether every stream compatible with the $J_3(\mathbb{O})$ ceiling saturates the bound, or only our stream saturates it, is not settled by the current apparatus. The Dixon / anomaly / Clifford perspectives each confirm saturation in our stream but do not establish that non-saturation is impossible in other physical-aspect basins.
  \item No phantom crossing --- structural in our stream (cuscuton kinetic positivity), but whether this is a generic-aspect topological theorem or a stream-specific consequence of the square-root choice depends on the uniqueness of the square-root kinetic form, which is itself the subject of the S7 structural prediction.
  \item 5D self-tuning as the appropriate dimension --- 5D is the minimal-dimension extension from 4D that admits the self-tuning mechanism, but it is not the top of the extension tower; Ruliad / string-landscape contexts imply higher-dimensional streams above. Whether our stream \emph{is} 5D or merely \emph{well-approximated} as 5D at accessible scales is a scope-class question the current apparatus cannot decide.
\end{itemize}

\paragraph{Reading the audit.}
The audit does not privilege any class over another.  Generic-physical-aspect quantities are structurally robust but shared with every compatible stream; this-stream quantities are empirically sharp but only identify our basin; saturation/ambiguous quantities are where the next round of work happens (either sharpening to generic or sharpening to this-stream).  The existence of all three classes, named at once, is the scope-controlled form of the monograph's map --- v2's honest answer to the question ``what kind of claim are you making?'' at every derived quantity.

%% ====================================================================
\section{A Structural Observation on the Radion Mass Scale}
\label{app:radion-thinness}
%% ====================================================================

\noindent\fbox{\parbox{\dimexpr\columnwidth-2\fboxsep-2\fboxrule\relax}{%
\textbf{Scope note.} This subsection is a \emph{structural observation}, not a derivation.  It reports a coincidence between two independently-derived scales and flags the coincidence as hypothesis-generating, under discipline.  Nothing in the main predictions of this monograph depends on the observation, and no empirical claim is made about what the coincidence would imply if examined further.
}}

\medskip

The monograph independently derives the radion mass (Chapter~\ref{ch:ncg}, P4 of Appendix~\ref{app:predictions}):
\begin{equation}
  m_\mathrm{rad} \approx 120~\mathrm{GeV}\,,
\end{equation}
with $99.7\%$ of the mass supplied by the NCG one-loop contribution and the remaining fraction by Goldberger--Wise backreaction.  This value sits within the uncertainty band of the electroweak scale and is, as noted in the conditional-caveat fbox for P4 (App.~\ref{app:pred-conditional}), robust as an order-of-magnitude claim while the precise numerical value is conditional on the GW-sector correction.

Separately, the Coherence Principle's basin/stream vocabulary (see Anchor, \emph{The Coherence Principle}, Library volume) treats coherent multi-scale systems as nested streams inside a larger structure.  The Principle does not itself specify what energy scale, if any, is associated with the transition between intra-stream dynamics and features of the wider structure.  It is a general-purpose claim about coherence maintenance across scales, not a specific physical prediction about any particular energy.

The observation flagged here is a \emph{coincidence of independently-derived scales}: Meridian's radion-mass derivation and the Anchor's basin/stream vocabulary arrive, from entirely independent apparatus (spectral action + GW stabilisation on the one hand, four coherence conditions on the other), at structures living at approximately the same scale (the electroweak window, $\sim 10^2~\mathrm{GeV}$).  This is observable; it is one sentence.  The hypothesis it generates --- that the radion-scale structure may be where this stream's intra-basin dynamics decouple from, or couple into, features that are properly generic-physical-aspect rather than this-stream --- is \emph{not} advanced here as a claim, derivation, or prediction.

We flag three scope-limits explicitly:
\begin{itemize}\setlength{\itemsep}{0pt}
  \item \textbf{No empirical claim about cross-basin accessibility.}  What it would mean, operationally, for a stream to be ``accessed'' from outside is a question the current apparatus cannot answer.  The Principle provides a vocabulary; it does not provide a cross-basin measurement procedure.  Nothing here claims that the electroweak scale is a portal to other basins.
  \item \textbf{No claim that the coincidence requires explanation.}  Independently-derived scales coincide at many points in physics (the Planck-scale / Higgs-scale hierarchy, the cosmological-constant / QCD-scale coincidence).  Coincidence is a flag for attention, not evidence of connection.  The observation here is reported at the same confidence as any other scale-coincidence in the monograph.
  \item \textbf{No claim drawn from outside the published record.}  The observation is derived from Meridian's own $m_\mathrm{rad}$ computation and the Anchor's own published basin/stream vocabulary.  No suppressed, non-public, or gatekept research is implied or cited.  The apparatus used is fully visible in the two volumes' bibliographies.
\end{itemize}

\paragraph{Why flag it at all.}
Because the Principle predicts that a coherent description should name its own scope-interesting features, and because an independently-derived coincidence between two volumes of the \emph{Corpus-Perspectival} library is a scope-interesting feature even in the absence of a derivation linking them.  Leaving the coincidence unreported would be a different failure of scope discipline: silence on a feature the Principle predicts should be noticed.  The compromise adopted here is to report the coincidence, fence it with discipline, and decline to advance any claim that would require apparatus this monograph does not possess.  Whether the coincidence proves meaningful, accidental, or misidentified is a question for later work in this or a higher framework.
